% =============================================================================
% Patch: jet-Gram foundation (referee-revised, hardened)
%   §2  sec:local-kernel-and-jet-normalization
%   §3  sec:jet-limit-local-blocks
%
% Target file:   rh/rh_rebuild.tex
% Target lines:  307 through 578 (§2 + §3, contiguous)
%
% Action: replace the existing §2 and §3 with the content below.
%
% Referee revisions integrated:
%   - theta defined precisely as Im L(1/4 + i t / 2) - (t/2) log pi, with
%     L(z) the unique holomorphic branch of log Gamma z on Re z > 0
%     that is real on the positive real axis (Definition
%     `def:riemann-siegel-phase').
%   - Window uniformity made explicit at section start:
%     I_T subset [T/2, 2T] and |I_T| << Q(T)^{-1}, used uniformly in
%     all Stirling-type estimates; differentiated Stirling on a fixed
%     sector recorded.
%   - Differentiated theta asymptotics extracted as an explicit lemma
%     (`lem:theta-derivative-asymptotics') giving uniform bounds on
%     q, q', q'' over I_T with absolute O-constants.
%   - (P2) recast as q(t) >= 1 (which subsumes any polynomial-in-Q
%     bound with c_q = 0 when Q >= 1).
%   - Diagonal kernel derivatives separated into their own lemma
%     (`lem:phase-kernel-diagonal-derivatives') with a *factored*
%     diagonal Taylor proof, so no remainder is divided by u - v.
%   - Same-point jet limit stated with explicit entrywise O(h^2) rate.
%   - Cross-block jet limit: hypothesis 0 < h < |s|/3 added; proof now
%     uses fourth-order parity argument with sign weights
%     1, sigma_2, sigma_1, sigma_1 sigma_2; the constraint h < |s|/3
%     is used explicitly to keep |s + u_1 - u_2| >= |s|/3 on the
%     sampled square.
%   - Same-point Gram positivity made explicitly theta-specific, with
%     two-sided lambda_min asymptotic 2 q(T) / pi (1 + o(1)).
%   - Scope guardrail remark added: positivity here does NOT prove
%     anisotropy lower bounds (e.g. |q'' - 6 q + 2 q^3|) or supply
%     source-energy information.
%   - Final hardening: theta remainder admits three differentiations
%     (q'' = theta'''), local phase chart marked as specialized
%     theta chart, diagonal-kernel proof made explicit at u = v,
%     P_h diagonal rescaling vs Gram-block congruence distinguished,
%     cross-block hypothesis includes neighborhood containment, error
%     constant depends on a finite power of |s|^{-1}, and a
%     fixed-separation scope remark prevents misuse as a collision
%     expansion.
%
% Closes:
%   §2: rem:wip-local-phase-chart
%   §3: rem:wip-same-point-jet-limit
%       rem:wip-cross-block-jet-limit
%       rem:wip-same-point-gram-positivity
%
% Status:
%   §2: \statusmig      not-started -> migrated
%       \statusreferee  not-started -> in-progress
%       \statussympy    not-started -> verified
%           [rh/sympy/sec-local-kernel-and-jet-normalization.py]
%       \statussim      not-started -> verified
%           [rh/simulation/sec-local-kernel-and-jet-normalization.py]
%       \statuslean     not-started -> in-progress
%           [rh/lean/RH/LocalKernelJetNormalization.lean]
%   §3: \statusmig      not-started -> migrated
%       \statusreferee  not-started -> in-progress
%       \statussympy    not-started -> verified
%           [rh/sympy/sec-jet-limit-local-blocks.py]
%       \statussim      not-started -> partial
%           [rh/simulation/sec-jet-limit-local-blocks.py]
%       \statuslean     not-started -> in-progress
%           [rh/lean/RH/JetLimitLocalBlocks.lean]
% =============================================================================

% =============================================================================
\section{Local kernel and jet normalization}
\label{sec:local-kernel-and-jet-normalization}
\statuspaper{LIVE}\quad
\statusmig{migrated}\quad
\statusreferee{in-progress}\quad
\statussympy[rh/sympy/sec-local-kernel-and-jet-normalization.py]{verified}\quad
\statussim[rh/simulation/sec-local-kernel-and-jet-normalization.py]{verified}\quad
\statuslean[rh/lean/RH/LocalKernelJetNormalization.lean]{in-progress}
\par\medskip

Fix a height parameter \(T\), and let \(I_T\) denote an interval
centered at \(T\) with
\[
        |I_T|\ll Q(T)^{-1},
        \qquad I_T\subset [T/2,2T].
\]
The Stirling estimates of
Lemma~\ref{lem:theta-derivative-asymptotics} are uniform on \(I_T\)
under this inclusion.

\subsection{The Riemann--Siegel phase chart}
\label{subsec:local-phase-chart}

\begin{definition}[Riemann--Siegel phase]
\label{def:riemann-siegel-phase}
Let \(L(z)\) be the unique holomorphic branch of \(\log\Gamma z\) on
\(\Re z>0\) for which \(L(x)=\log\Gamma(x)\in\mathbb R\) when \(x>0\).
For \(t>0\), define
\begin{equation}
\label{eq:theta-definition}
 \theta(t)
 := \Im L\!\left(\frac14+\frac{it}{2}\right)
    -\frac{t}{2}\log\pi,
\end{equation}
and \(Z(t):=e^{i\theta(t)}\zeta(1/2+it)\).  The functional equation
gives
\[
        \zeta\!\left(\tfrac12+it\right)
        =e^{-i\theta(t)}\,Z(t),
        \qquad Z(t)\in\mathbb R \text{ for } t>0.
\]
\end{definition}

\begin{lemma}[Differentiated theta asymptotics]
\label{lem:theta-derivative-asymptotics}
Uniformly for \(t\in I_T\),
\begin{align}
\label{eq:theta-derivative-asymptotics}
q(t):=\theta'(t)
&=\frac12\log\frac{t}{2\pi}-\frac{1}{48 t^2}+O(t^{-4}),\\
q'(t)=\theta''(t)
&=\frac{1}{2t}+O(t^{-3}),\\
q''(t)=\theta'''(t)
&=-\frac{1}{2t^2}+O(t^{-4}).
\end{align}
The constants in the \(O\)-terms are absolute.
\end{lemma}

\begin{proof}
Let \(z=1/4+it/2\).  The half-plane \(\Re z\ge 1/4\) lies in a fixed
Stirling sector, and Stirling's expansion for \(\log\Gamma z\), with
remainder differentiated a fixed number of times, is uniform there as
\(|z|\to\infty\).  Taking imaginary parts in \eqref{eq:theta-definition}
gives the standard Riemann--Siegel expansion
\[
\theta(t)=\frac{t}{2}\log\frac{t}{2\pi}-\frac{t}{2}-\frac{\pi}{8}
          +\frac{1}{48 t}+O(t^{-3}),
\]
again with the displayed remainder admitting three differentiations
with the corresponding differentiated bounds.  Differentiating gives
\(q(t)=\theta'(t)\) as stated; differentiating once and twice more gives
the formulas for \(q'(t)\) and \(q''(t)\).  Since
\(I_T\subset[T/2,2T]\), these estimates are uniform on \(I_T\), with
absolute constants.
\end{proof}

\begin{definition}[Local phase chart]
\label{def:local-phase-chart}
A retained high-height local phase chart is the interval \(I_T\)
together with the real smooth phase
\[
        \Ph_T(t):=\theta(t)\qquad (t\in I_T).
\]
When \(T\) is fixed by context, \(\Ph := \Ph_T\) and \(q := \Ph'\).
This is a specialized theta chart, not a general phase-chart axiom.
The chart conditions used downstream are:
\begin{enumerate}[label=\textup{(P\arabic*)},leftmargin=*]
\item \(\Ph_T = \theta\) on \(I_T\);
\item \(q(t)\ge 1\) on \(I_T\), after increasing the global lower-height
      cutoff if necessary.
\end{enumerate}
If \(Q(T)\ge 1\), then (P2) gives the polynomial-in-\(Q\) lower bound
\(q(t)\ge Q(T)^{-c_q}\) with \(c_q = 0\).  The uniform estimates
\(q(t)\asymp \log T\), \(q'(t)=O(T^{-1})\), and
\(q''(t)=O(T^{-2})\) on \(I_T\) are supplied by
Lemma~\ref{lem:theta-derivative-asymptotics}.
\end{definition}

\begin{lemma}[Phase-derivative lower bound]
\label{lem:phase-derivative-lower-bound}
Condition \textup{(P2)} of Definition~\ref{def:local-phase-chart}
holds for all sufficiently large \(T\).  More quantitatively, uniformly
on \(I_T\),
\[
        q(t)\ge \frac12\log\frac{T}{4\pi}-C\,T^{-2}
\]
for an absolute constant \(C\).
\end{lemma}

\begin{proof}
By Lemma~\ref{lem:theta-derivative-asymptotics},
\(q(t)=\tfrac12\log(t/(2\pi))+O(t^{-2})\).  Since
\(t \in I_T \subset [T/2, 2T]\),
\(\log(t/(2\pi)) \ge \log(T/(4\pi))\), and the stated bound follows.
The right-hand side is eventually larger than \(1\).
\end{proof}

\begin{remark}[No \(\zeta\)-side regularity is used here]
\label{rem:chart-regularity-zeta-side}
The chart uses \(\zeta\) only through the identity
\(\zeta(1/2+it)=e^{-i\theta(t)}Z(t)\) and through the theta phase.
No lower bound on \(|Z|\), no zero-free condition in \(I_T\), and no
source-energy input is used in
Sections~\ref{sec:local-kernel-and-jet-normalization}--\ref{sec:jet-limit-local-blocks}.
Such conditions are imposed only in the later sections where
actual-zeta transverse suppression or source-energy transfer is
consumed.
\end{remark}

\subsection{The phase kernel}
\label{subsec:phase-kernel}

\begin{definition}[Phase kernel]
\label{def:phase-kernel}
For a real \(C^1\) phase \(\Ph\) on an interval \(I\),
\begin{equation}
\label{eq:phase-kernel}
K_\Ph(x,y)=
\begin{cases}
\dfrac{\sin\bigl(\Ph(x)-\Ph(y)\bigr)}{\pi(x-y)}, & x\ne y,\\[1.2ex]
\dfrac{q(x)}{\pi}, & x=y,
\end{cases}
\qquad q=\Ph'.
\end{equation}
\end{definition}

\begin{lemma}[Basic kernel properties]
\label{lem:phase-kernel-properties}
\(K_\Ph\) is real-valued and symmetric.  If \(\Ph\in C^1\), the
extension across the diagonal is continuous and
\[
        K_\Ph(T,T)=\frac{q(T)}{\pi}.
\]
If \(\Ph\in C^4\), the diagonal first and mixed second derivatives
exist and are given by
Lemma~\ref{lem:phase-kernel-diagonal-derivatives}.
\end{lemma}

\begin{proof}
For \(x\ne y\), symmetry follows from
\(\sin(\Ph(y)-\Ph(x))/(y-x) = \sin(\Ph(x)-\Ph(y))/(x-y)\).
Real-valuedness is immediate.  The diagonal value follows from
\(\Ph(y+h)-\Ph(y) = h q(y) + o(h)\).
\end{proof}

\begin{lemma}[Off-diagonal kernel derivatives]
\label{lem:phase-kernel-derivatives}
For \(T_1\ne T_2\) with \(s=T_1-T_2\), \(\Del=\Ph(T_1)-\Ph(T_2)\),
\(q_j=q(T_j)\), and \(\Ph\in C^1\) near \(T_1, T_2\),
\begin{align*}
K_\Ph(T_1,T_2) &= \frac{\sin\Del}{\pi s},
&K_x(T_1,T_2) &= \frac{q_1\,s\cos\Del - \sin\Del}{\pi s^2},\\
K_y(T_1,T_2) &= \frac{\sin\Del - q_2\,s\cos\Del}{\pi s^2},
&K_{xy}(T_1,T_2)
&= \frac{(q_1+q_2)\,s\cos\Del + (q_1 q_2 s^2 - 2)\sin\Del}{\pi s^3}.
\end{align*}
\end{lemma}

\begin{proof}
Differentiate \(K_\Ph(x,y) = \sin(\Ph(x)-\Ph(y))/\pi(x-y)\) by the
quotient and chain rules, then evaluate at \((T_1, T_2)\).
\end{proof}

\begin{lemma}[Diagonal kernel derivatives]
\label{lem:phase-kernel-diagonal-derivatives}
If \(\Ph\in C^4\) near \(T\), then
\begin{equation}
\label{eq:diagonal-kernel-derivatives}
K_\Ph(T,T)=\frac{q(T)}{\pi},\qquad
K_x(T,T)=K_y(T,T)=\frac{q'(T)}{2\pi},
\end{equation}
\[
        K_{xy}(T,T)=\frac{q''(T)+2q(T)^3}{6\pi}.
\]
\end{lemma}

\begin{proof}
Set \(x=T+u\), \(y=T+v\), and write \(r=|u|+|v|\); derivatives below
are evaluated at \(T\).  Factor the phase difference as
\[
\Ph(T+u)-\Ph(T+v)=(u-v)A(u,v),\qquad
A(u,v)=\int_0^1 q\bigl(T+v+\tau(u-v)\bigr)\,d\tau,
\]
where the integral form extends continuously to \(u=v\).  Taylor
expansion of the integrand gives
\[
A(u,v)
= q+\frac{q'}2(u+v)+\frac{q''}{6}(u^2+uv+v^2)+O(r^3).
\]
For \(u\ne v\), and then by continuous extension to \(u=v\),
\(\Ph(T+u)-\Ph(T+v)=(u-v)A(u,v)=O(r)\) gives
\[
\frac{\sin(\Ph(T+u)-\Ph(T+v))}{u-v}
= A(u,v)-\frac{(u-v)^2}{6}A(u,v)^3+O(r^3).
\]
Keeping terms through total degree two gives
\[
K_\Ph(T+u,T+v)
= \frac1\pi\!\left[
 q + \frac{q'}{2}(u+v)
 + \!\left(\frac{q''}{6} - \frac{q^3}{6}\right)\!(u^2+v^2)
 + \!\left(\frac{q''}{6} + \frac{q^3}{3}\right)\!uv
 + O(r^3)\right].
\]
The coefficients of \(u\), \(v\), and \(uv\) are respectively
\(K_x(T,T)\), \(K_y(T,T)\), and \(K_{xy}(T,T)\), proving the claim.
\end{proof}

\subsection{Point-to-jet transform}
\label{subsec:point-to-jet-transform}

For symmetric pairs \(T\pm h\), define the Gram-normalized coalescing
even/odd transform
\begin{equation}
\label{eq:point-to-jet-transform}
P_h
=\frac1{\sqrt2}
\begin{pmatrix}
1 & 1\\[0.6ex]
-1/(2h) & 1/(2h)
\end{pmatrix}
=\diag(\sqrt 2, 1/\sqrt 2)\,
\begin{pmatrix}
1/2 & 1/2\\[0.6ex]
-1/(2h) & 1/(2h)
\end{pmatrix}.
\end{equation}
Thus \(P_h\) is the literal central value/derivative jet transform
followed by the fixed diagonal rescaling
\(\diag(\sqrt 2, 1/\sqrt 2)\); at the level of Gram blocks this induces
the corresponding diagonal congruence.  The value/derivative basis
produced by \(P_h\) is the \emph{Gram-normalized basis}.

With this convention,
\begin{equation}
\label{eq:same-point-as-kernel-derivatives}
J(T)=
\begin{pmatrix}
2K_\Ph(T,T) & \partial_y K_\Ph(T,T)\\[0.6ex]
\partial_x K_\Ph(T,T) & \tfrac12 \partial_{xy} K_\Ph(T,T)
\end{pmatrix},
\end{equation}
which is the same-point block computed explicitly in
Lemma~\ref{lem:same-point-jet-limit}.

\begin{remark}[Pair-separation convention]
\label{rem:pair-separation-convention}
The pair is \(T-h, T+h\), with separation \(2h\); the lower row of
\(P_h\) carries \(1/(2h)\).  Any formula imported from the alternative
\(T\pm h/2\) convention or from the literal-jet matrix with lower row
\((-1/h, 1/h)\) must be re-conjugated before use.
\end{remark}

\begin{remark}[Convention check]
\label{rem:normalization-audit}
With the literal central value/derivative transform at the pair
\(T\pm h\), the same-point limit would be
\[
\frac1\pi
\begin{pmatrix}
q(T) & \tfrac12 q'(T)\\[0.6ex]
\tfrac12 q'(T) & \tfrac16\bigl(q''(T)+2q(T)^3\bigr)
\end{pmatrix},
\]
related to \(J(T)\) by the diagonal congruence
\(\diag(\sqrt 2, 1/\sqrt 2)\) of \eqref{eq:point-to-jet-transform}; the
\((1,1)\) entry doubles and the \((2,2)\) entry is halved.  All
same-point blocks, cross-blocks, finite-scale blocks, whitening
matrices, and quotient covectors in this paper are read in the
Gram-normalized convention.
\end{remark}

% =============================================================================
\section{Jet-limit local blocks}
\label{sec:jet-limit-local-blocks}
\statuspaper{LIVE}\quad
\statusmig{migrated}\quad
\statusreferee{in-progress}\quad
\statussympy[rh/sympy/sec-jet-limit-local-blocks.py]{verified}\quad
\statussim[rh/simulation/sec-jet-limit-local-blocks.py]{partial}\quad
\statuslean[rh/lean/RH/JetLimitLocalBlocks.lean]{in-progress}
\par\medskip

\subsection{Same-point jet limit}
\label{subsec:same-point-jet-limit}

\begin{lemma}[Same-point jet limit]
\label{lem:same-point-jet-limit}
Assume \(\Ph\in C^5\) near \(T\).  Let
\[
A_h(T)=
\begin{pmatrix}
K_\Ph(T-h,T-h) & K_\Ph(T-h,T+h)\\
K_\Ph(T+h,T-h) & K_\Ph(T+h,T+h)
\end{pmatrix}
\]
with row and column order \((T-h, T+h)\).  Then
\[
        P_h\,A_h(T)\,P_h^\top \;=\; J(T) + O(h^2)
        \qquad (h\to 0),
\]
where the \(O(h^2)\) is entrywise, with a constant depending only on a
fixed neighborhood of \(T\) and on a local \(C^5\)-bound for \(\Ph\).  The
limiting block is
\begin{equation}
\label{eq:same-point-J}
J(T)=\frac1\pi
\begin{pmatrix}
2q(T) & \tfrac12 q'(T)\\[1ex]
\tfrac12 q'(T) & \tfrac1{12}\bigl(q''(T)+2q(T)^3\bigr)
\end{pmatrix}.
\end{equation}
\end{lemma}

\begin{proof}
Derivatives below are evaluated at \(T\).  Taylor expansion gives
\[
\Ph(T\pm h) = \Ph(T) \pm hq + \frac{h^2}{2}q'
            \pm \frac{h^3}{6}q'' + \frac{h^4}{24}q'''
            + O(h^5),
\]
so
\[
\Ph(T-h)-\Ph(T+h) = -2hq - \frac{h^3}{3}q'' + O(h^5).
\]
Using \(\sin z = z - z^3/6 + O(z^5)\),
\begin{equation}
\label{eq:A12-same-point-expansion}
K_\Ph(T-h, T+h)
= \frac{q}{\pi} + \frac{h^2}{6\pi}(q'' - 4q^3) + O(h^4).
\end{equation}
Also \(K_\Ph(T\pm h,T\pm h)=q(T\pm h)/\pi\), with
\[
q(T\pm h)=q\pm hq'+\frac{h^2}{2}q''
          \pm\frac{h^3}{6}q'''+O(h^4).
\]
For symmetric \(M\), conjugation by \eqref{eq:point-to-jet-transform}
gives
\begin{align*}
(P_h M P_h^\top)_{11} &= \tfrac12(M_{11}+2M_{12}+M_{22}),\\
(P_h M P_h^\top)_{12} &= \tfrac{1}{4h}(M_{22}-M_{11}),\\
(P_h M P_h^\top)_{22} &= \tfrac{1}{8h^2}(M_{11}-2M_{12}+M_{22}).
\end{align*}
Substituting \eqref{eq:A12-same-point-expansion} and the Taylor series
of \(q(T\pm h)\),
\begin{align*}
(P_h A_h P_h^\top)_{11} &= \frac{2q}{\pi} + O(h^2),\\
(P_h A_h P_h^\top)_{12} &= \frac{q'}{2\pi} + O(h^2),\\
(P_h A_h P_h^\top)_{22} &= \frac{q''+2q^3}{12\pi} + O(h^2),
\end{align*}
which is \eqref{eq:same-point-J}.
\end{proof}

\subsection{Cross-block jet limit}
\label{subsec:cross-block-jet-limit}

\begin{lemma}[Cross-block jet limit]
\label{lem:cross-block-jet-limit}
Assume \(\Ph\in C^5\) in neighborhoods of \(T_1\) and \(T_2\), let
\(T_1\ne T_2\), set
\[
        s = T_1 - T_2,
        \qquad \Del = \Ph(T_1)-\Ph(T_2),
        \qquad q_j = q(T_j),
\]
and assume \(0 < h < |s|/3\), with \(h\) also sufficiently small that
the four sampled points lie in the chosen neighborhoods of \(T_1\) and
\(T_2\).  Define
\[
C_h(T_1, T_2)=
\begin{pmatrix}
K_\Ph(T_1-h, T_2-h) & K_\Ph(T_1-h, T_2+h)\\
K_\Ph(T_1+h, T_2-h) & K_\Ph(T_1+h, T_2+h)
\end{pmatrix}
\]
with row order \((T_1-h, T_1+h)\) and column order \((T_2-h, T_2+h)\).
Then
\[
        P_h\,C_h(T_1, T_2)\,P_h^\top
        \;=\; \frac1\pi N_{12}(T_1, T_2) + O(h^2)
        \qquad (h\to 0),
\]
where the \(O(h^2)\) is entrywise for fixed \(T_1\ne T_2\); the implied
constant may depend on a finite power of \(|s|^{-1}\) and on local
\(C^5\)-bounds for \(\Ph\), but not on \(h\).  Here
\begin{equation}
\label{eq:cross-N12}
N_{12}=
\begin{pmatrix}
\dfrac{2\sin\Del}{s}
&
\dfrac{\sin\Del - q_2 s\cos\Del}{s^2}
\\[2.4ex]
\dfrac{q_1 s\cos\Del - \sin\Del}{s^2}
&
\dfrac{(q_1+q_2)s\cos\Del + (q_1 q_2 s^2 - 2)\sin\Del}{2s^3}
\end{pmatrix}.
\end{equation}
\end{lemma}

\begin{proof}
The kernel is smooth in a neighborhood of \((T_1, T_2)\), and the
hypothesis \(h<|s|/3\) gives
\(|s+u_1-u_2|\ge |s|/3\) for \(|u_1|,|u_2|\le h\), bounding the
Taylor remainder on the sampled square.  Write
\(K_{ab} = \partial_x^a \partial_y^b K_\Ph(T_1, T_2)\).  Taylor
expansion to total order four gives, uniformly for
\(u_1, u_2 \in \{-h, +h\}\),
\[
K_\Ph(T_1+u_1, T_2+u_2)
= \sum_{a+b\le 4}\frac{K_{ab}}{a!\,b!}\, u_1^a\, u_2^b + O_s(h^5).
\]
The subscript indicates that the constant may depend on a finite power
of \(|s|^{-1}\) and on the chosen off-diagonal neighborhoods.  For a
general \(2\times 2\) matrix \(M\),
\begin{align*}
(P_h M P_h^\top)_{11}
&= \tfrac12(M_{11} + M_{12} + M_{21} + M_{22}),\\
(P_h M P_h^\top)_{12}
&= \tfrac{1}{4h}(-M_{11} + M_{12} - M_{21} + M_{22}),\\
(P_h M P_h^\top)_{21}
&= \tfrac{1}{4h}(-M_{11} - M_{12} + M_{21} + M_{22}),\\
(P_h M P_h^\top)_{22}
&= \tfrac{1}{8h^2}(M_{11} - M_{12} - M_{21} + M_{22}).
\end{align*}
Equivalently, these four expressions apply the sign weights
\(1\), \(\sigma_2\), \(\sigma_1\), \(\sigma_1\sigma_2\) respectively to
the four samples with \(u_i = \sigma_i h\).  The parity of the Taylor
monomials therefore selects
\[
(P_h C_h P_h^\top)_{11} = 2 K_{00} + O(h^2),\qquad
(P_h C_h P_h^\top)_{12} = K_{01} + O(h^2),
\]
\[
(P_h C_h P_h^\top)_{21} = K_{10} + O(h^2),\qquad
(P_h C_h P_h^\top)_{22} = \tfrac12 K_{11} + O(h^2).
\]
Substituting the off-diagonal derivative formulas of
Lemma~\ref{lem:phase-kernel-derivatives} gives \eqref{eq:cross-N12}.
\end{proof}

\begin{remark}[Fixed-separation scope]
\label{rem:cross-block-fixed-separation-scope}
Lemma~\ref{lem:cross-block-jet-limit} is an off-diagonal jet limit at
fixed separation.  It is not a collision expansion for \(T_1\to T_2\);
any two-center blow-up or near-coalescence analysis must be derived
from a separate scaling argument.
\end{remark}

\subsection{Same-point Gram positivity}
\label{subsec:same-point-gram-positivity}

\begin{definition}[Same-point Gram determinant]
\label{def:same-point-gram-determinant}
\begin{equation}
\label{eq:same-point-gram-determinant}
D_J(T) \;:=\; 4\,q(T)^4 + 2\,q(T)\,q''(T) - 3\,q'(T)^2.
\end{equation}
\end{definition}

\begin{lemma}[Algebraic same-point Gram criterion]
\label{lem:algebraic-same-point-gram-criterion}
For \(J(T)\) of \eqref{eq:same-point-J},
\begin{equation}
\label{eq:J-trace-det}
\Tr J(T) = \frac{2 q(T)^3 + 24\,q(T) + q''(T)}{12\pi},\qquad
\det J(T) = \frac{D_J(T)}{12\pi^2}.
\end{equation}
Moreover,
\[
        J(T)\succ 0
        \quad\Longleftrightarrow\quad
        q(T) > 0 \quad\text{and}\quad D_J(T) > 0.
\]
\end{lemma}

\begin{proof}
The trace and determinant follow by direct computation from
\eqref{eq:same-point-J}.  Since \(J(T)\) is real symmetric,
Sylvester's criterion gives positive definiteness exactly when its
leading principal minor \(2q(T)/\pi\) and its determinant are
positive.
\end{proof}

\begin{lemma}[Same-point Gram positivity for the theta phase]
\label{lem:same-point-gram-positivity}
For \(\Ph_T = \theta\), and for all sufficiently large retained
heights, \(J(T)\) is positive definite.  More precisely,
\begin{equation}
\label{eq:lambda-min-asymptotic}
        \lambda_{\min}(J(T))
        = \frac{2 q(T)}{\pi}\,(1 + o(1))
        = \frac{1}{\pi}\,\log\frac{T}{2\pi}\,(1 + o(1)).
\end{equation}
After increasing the lower-height cutoff if necessary,
\(\lambda_{\min}(J(T)) \ge 1\); hence the polynomial-in-\(Q\) lower
bound \(\lambda_{\min}(J(T)) \ge Q(T)^{-C_J}\) holds with \(C_J = 0\)
whenever \(Q(T) \ge 1\).
\end{lemma}

\begin{proof}
By Lemma~\ref{lem:theta-derivative-asymptotics},
\[
q(T) = \tfrac12 \log(T/(2\pi)) + O(T^{-2}),\qquad
q'(T) = O(T^{-1}),\qquad q''(T) = O(T^{-2}).
\]
Therefore
\[
D_J(T) = 4\,q(T)^4 + O\!\bigl(q(T)\,T^{-2}\bigr) + O(T^{-2})
       = 4\,q(T)^4\,(1 + o(1)) > 0
\]
for sufficiently large \(T\), and \(q(T) > 0\).  By
Lemma~\ref{lem:algebraic-same-point-gram-criterion},
\(J(T) \succ 0\).

Let \(0 < \lambda_1 \le \lambda_2\) be the eigenvalues of \(J(T)\).
Since \(\lambda_2 \le \Tr J(T)\),
\[
\lambda_1 = \frac{\det J(T)}{\lambda_2}
          \ge \frac{\det J(T)}{\Tr J(T)}
          = \frac{D_J(T)}{\pi(2 q(T)^3 + 24\,q(T) + q''(T))}
          = \frac{2 q(T)}{\pi}\,(1 + o(1)).
\]
The Rayleigh quotient at \(e_1\) gives
\(\lambda_1 \le J_{11}(T) = 2 q(T)/\pi\), proving
\eqref{eq:lambda-min-asymptotic}.
\end{proof}

\begin{remark}[Scope of the positivity lemma]
\label{rem:scope-of-gram-positivity}
Lemma~\ref{lem:same-point-gram-positivity} is theta-specific.  For an
arbitrary real phase, the safe statement is only the algebraic
criterion \(q > 0\) and \(D_J > 0\).  This lemma does not prove any
lower bound for unrelated anisotropy functionals such as
\(q'' - 6q + 2q^3\), nor does it supply source-energy input of the
kind consumed in Section~\ref{sec:source-energy-supply}.
\end{remark}
